% !TEX root = ../main.tex
\corrigesChapitre{10}{Électrolyse et stockage de l'énergie}
\corrigesSouples

\begin{blocLivre}{cor-10-1}
\begin{corrige}{1}{Pile ou électrolyseur ?}
\demarche{Spontané : Qr se rapproche de K ; forcé : le générateur impose le sens inverse et Qr s'éloigne de K.}
\begin{questions}
\item Dans le montage A : le système fonctionne en pile et évolue spontanément selon
$\res{\ce{Cu^2+(aq) + Zn(s) -> Cu(s) + Zn^2+(aq)}}$.
\item Montage A : $Q_\mathrm{r}$ \res{\text{se rapproche de }K} ; montage B : le courant inverse fait croître
$[\ce{Cu^2+}]$ et décroître $[\ce{Zn^2+}]$, $Q_\mathrm{r}$ \res{\text{s'éloigne de }K}.
\item $\res{\ce{Cu(s) + Zn^2+(aq) -> Cu^2+(aq) + Zn(s)}}$ : ce sens est contraire au critère d'évolution, il n'est
possible que grâce à l'énergie apportée par le générateur, d'où le nom de transformation forcée.
\item L'anode est le siège de l'oxydation : montage A, \res{\text{lame de zinc}} ; montage B, \res{\text{lame de
cuivre}}, reliée à la borne $+$ du générateur.
\end{questions}
\end{corrige}
\end{blocLivre}
\begin{blocLivre}{cor-10-2}
\begin{corrige}{2}{Les électrodes d'un électrolyseur}
\demarche{Le dépôt rougeâtre est du cuivre (réduction) ; la coloration orangée est due au dibrome (oxydation).}
\begin{questions}
\item Le dépôt rougeâtre est du \res{\text{cuivre}} ; la coloration orangée révèle du \res{\text{dibrome}} dissous.
\item Électrode 1 : \ce{Cu^2+ + 2 e- -> Cu} ; électrode 2 : \ce{2 Br- -> Br2 + 2 e-}.
\item L'électrode 1, siège d'une réduction, est la \res{\text{cathode}}, reliée à la borne $-$ ; l'électrode 2, siège
d'une oxydation, est l'\res{\text{anode}}, reliée à la borne $+$.
\item $\res{\ce{Cu^2+(aq) + 2 Br^-(aq) -> Cu(s) + Br2(aq)}}$. Sans générateur, plus rien ne force la transformation :
l'électrolyse s'arrête et le système, devenu une pile, tendrait à évoluer en sens inverse.
\end{questions}
\end{corrige}
\end{blocLivre}
\begin{blocLivre}{cor-10-3}
\begin{corrige}{3}{Quantité d'électricité et quantité de matière}
\demarche{$Q=I\,\Delta t$, puis $n(\mathrm{e}^-)=Q/F$ ; la demi-équation relie $n(\mathrm{e}^-)$ à la quantité d'argent.}
\begin{questions}
\item $Q=I\,\Delta t=0{,}50\times1200=\res{\qty{600}{C}}$.
\item $n(\mathrm{e}^-)=\dfrac{600}{\num{96500}}=\res{\qty{6,22e-3}{mol}}$, soit
$\num{6.22e-3}\times\num{6.02e23}=\res{\num{3,74e21}}$ électrons.
\item Un électron par ion argent : $n(\ce{Ag})=\res{\qty{6,22e-3}{mol}}$.
\item La quantité déposée est proportionnelle à la durée : il faut \res{\qty{40}{min}}.
\end{questions}
\end{corrige}
\end{blocLivre}
\begin{blocLivre}{cor-10-4}
\begin{corrige}{4}{Cuivrer une médaille}
\demarche{$n(\ce{Cu})=n(\mathrm{e}^-)/2$, puis masse, volume ($m/\rho$) et épaisseur ($V/S$).}
\begin{questions}
\item \ce{Cu^2+ + 2 e- -> Cu} : c'est une réduction, elle se produit à la \res{\text{cathode}}, la médaille étant
reliée à la borne $-$.
\item $n(\mathrm{e}^-)=\dfrac{0{,}60\times2400}{\num{96500}}=\qty{1,49e-2}{mol}$, donc
$n(\ce{Cu})=\qty{7,46e-3}{mol}$ et $m=\num{7.46e-3}\times63{,}5=\res{\qty{0,47}{g}}$.
\item $V=\dfrac{m}{\rho}=\dfrac{0{,}474}{8{,}96}=\qty{5,3e-2}{cm^3}$, d'où
$e=\dfrac{V}{S}=\dfrac{\num{5.3e-2}}{60}=\qty{8,8e-4}{cm}=\res{\qty{8,8}{\micro m}}$.
\item La masse déposée est proportionnelle à la durée : \res{\text{doubler la durée}} (80~min).
\end{questions}
\end{corrige}
\end{blocLivre}
\begin{blocLivre}{cor-10-5}
\begin{corrige}{5}{Recharger un accumulateur}
\demarche{En charge, chaque demi-équation s'inverse et les rôles anode/cathode s'échangent.}
\begin{questions}
\item Décharge : au zinc (pôle $-$), \ce{Zn -> Zn^2+ + 2 e-} ; au cuivre (pôle $+$), \ce{Cu^2+ + 2 e- -> Cu}.
\item Charge : au cuivre, \ce{Cu -> Cu^2+ + 2 e-} : le dépôt \res{\text{se dissout}} ; au zinc,
\ce{Zn^2+ + 2 e- -> Zn} : du zinc se dépose.
\item Pendant la charge, l'oxydation a lieu sur le cuivre : \res{\text{l'électrode de cuivre est l'anode}}, alors
qu'elle était la cathode pendant la décharge.
\item Le système évolue à l'opposé du sens spontané ($Q_\mathrm{r}$ s'éloigne de $K$) : seul l'apport d'énergie
électrique du générateur le permet. Cette énergie est \res{\text{stockée sous forme chimique}} par les réactifs
reconstitués.
\end{questions}
\end{corrige}
\end{blocLivre}
\begin{blocLivre}{cor-10-6}
\begin{corrige}{6}{L'électrolyse de l'eau}
\demarche{Deux électrons par molécule de dihydrogène, quatre par molécule de dioxygène : les volumes sont dans le rapport 2.}
\begin{questions}
\item Cathode : \ce{2 H2O + 2 e- -> H2 + 2 HO-} ; anode : \ce{4 HO- -> O2 + 2 H2O + 4 e-} ; bilan :
$\res{\ce{2 H2O -> 2 H2 + O2}}$.
\item $n(\mathrm{e}^-)=\dfrac{0{,}25\times900}{\num{96500}}=\res{\qty{2,33e-3}{mol}}$.
\item $n(\ce{H2})=\dfrac{n(\mathrm{e}^-)}{2}=\qty{1,17e-3}{mol}$, soit $V(\ce{H2})=\res{\qty{28}{mL}}$ ;
$n(\ce{O2})=\dfrac{n(\mathrm{e}^-)}{4}$, soit $V(\ce{O2})=\res{\qty{14}{mL}}$ : le volume de dihydrogène est
\res{\text{le double}} de celui de dioxygène.
\item Le dihydrogène produit une petite détonation à l'approche d'une flamme ; le dioxygène rallume une bûchette
incandescente.
\end{questions}
\schema[corrige]{10-ex06}
\end{corrige}
\end{blocLivre}
\begin{blocLivre}{cor-10-7}
\begin{corrige}{7}{Argenter des couverts}
\demarche{Un électron par atome d'argent ; l'anode soluble compense exactement les ions consommés à la cathode.}
\begin{questions}
\item \ce{Ag+ + e- -> Ag}.
\item $n(\mathrm{e}^-)=n(\ce{Ag})=\dfrac{2{,}0}{107{,}9}=\qty{1,85e-2}{mol}$, soit $Q=\qty{1,79e3}{C}$ et
$\Delta t=\dfrac{Q}{I}=\dfrac{1789}{0{,}50}=\qty{3,6e3}{s}$, soit \res{\text{environ 60~min}}.
\item \ce{Ag -> Ag+ + e-}.
\item À chaque instant, l'anode libère autant d'ions \ce{Ag+} que la cathode en consomme : la concentration du bain
\res{\text{reste constante}}. Le dépôt garde ainsi la même vitesse et la même qualité pendant toute l'opération.
\end{questions}
\end{corrige}
\end{blocLivre}
\begin{blocLivre}{cor-10-8}
\begin{corrige}{8}{L'affinage du cuivre}
\demarche{Anode soluble en cuivre impur, cathode en cuivre pur ; l'énergie vaut E = U Q.}
\begin{questions}
\item Anode : \ce{Cu -> Cu^2+ + 2 e-} ; cathode : \ce{Cu^2+ + 2 e- -> Cu}. Ce qui est dissous à l'anode est déposé à
la cathode : $[\ce{Cu^2+}]$ \res{\text{ne varie pas}}.
\item $n(\mathrm{e}^-)=\dfrac{250\times86400}{\num{96500}}=\qty{224}{mol}$, donc $n(\ce{Cu})=\qty{112}{mol}$ et
$m=112\times63{,}5=\res{\qty{7,1}{kg}}$.
\item $E=U\,Q=0{,}30\times\num{2.16e7}=\qty{6,5e6}{J}=\res{\qty{1,8}{kW.h}}$, soit
$\dfrac{1{,}8}{7{,}1}=\res{\qty{0,25}{kW.h}}$ par kilogramme.
\item Dans l'eau, les ions \ce{Na+} ne peuvent pas être réduits avant l'eau elle-même : à la cathode, c'est l'eau qui
serait réduite en dihydrogène. L'électrolyse en solution aqueuse ne peut pas déposer un métal \res{\text{plus réducteur
que l'eau}}.
\end{questions}
\end{corrige}
\end{blocLivre}
\begin{blocLivre}{cor-10-9}
\begin{corrige}{9}{L'accumulateur au plomb}
\demarche{Capacité en coulombs, puis deux électrons par atome de plomb oxydé.}
\begin{questions}
\item Le plomb est oxydé : électrode de plomb $=$ \res{\text{anode}} (pôle $-$) ; l'oxyde \ce{PbO2} est réduit :
\res{\text{cathode}} (pôle $+$). Bilan :
$\res{\ce{Pb + PbO2 + 4 H+ + 2 SO4^2- -> 2 PbSO4 + 2 H2O}}$.
\item $Q=60\times3600=\qty{2,16e5}{C}$ et $n(\mathrm{e}^-)=\dfrac{\num{2.16e5}}{\num{96500}}=\res{\qty{2,24}{mol}}$.
\item $n(\ce{Pb})=\dfrac{n(\mathrm{e}^-)}{2}=\qty{1,12}{mol}$, soit $m=1{,}12\times207{,}2=\res{\qty{232}{g}}$.
\item $Q_\text{dém}=150\times2{,}0=\qty{300}{C}$, soit $\dfrac{300}{\num{2.16e5}}=\res{0{,}14~\%}$ de la capacité.
\item Les réactions s'inversent : à l'électrode de plomb, \ce{PbSO4 + 2 e- -> Pb + SO4^2-} ; à l'autre,
\ce{PbSO4 + 2 H2O -> PbO2 + 4 H+ + SO4^2- + 2 e-}.
\end{questions}
\end{corrige}
\end{blocLivre}
\begin{blocLivre}{cor-10-10}
\begin{corrige}{10}{Produire du dihydrogène}
\demarche{Deux électrons par molécule de dihydrogène : m = M I t/(2F).}
\begin{questions}
\item Pendant $t$, la quantité d'électrons vaut $I\,t/F$ et $n(\ce{H2})=I\,t/(2F)$ : la masse vaut
$m=M\,I\,t/(2F)$, ce qui est la ligne du programme.
\item Le programme avance heure par heure : au bout de \res{\qty{64800}{s}}, soit \res{18~h}, la masse atteint
\res{\qty{1007}{g}}, première valeur qui dépasse 1,0~kg.
\item $t=\dfrac{2F\,m}{M\,I}=\dfrac{2\times\num{96500}\times1000}{2{,}0\times1500}=\qty{6,43e4}{s}$, soit
\res{\qty{17,9}{h}} : le programme arrondit à l'heure entière suivante.
\item $n(\ce{H2})=500$ mol, donc $V=500\times24{,}8=\qty{1,24e4}{L}=\res{\qty{12,4}{m^3}}$ : un volume énorme, que
l'on réduit en comprimant le gaz à plusieurs centaines de bars ou en le liquéfiant.
\end{questions}
\end{corrige}
\end{blocLivre}
\begin{blocLivre}{cor-10-11}
\begin{corrige}{11}{Protéger l'acier par électrozingage}
\demarche{Volume = surface × épaisseur, puis masse, quantité d'électrons et durée.}
\begin{questions}
\item \ce{Zn^2+ + 2 e- -> Zn}.
\item $V=S\,e=0{,}96\times\num{10e-6}=\qty{9,6e-6}{m^3}=\qty{9,6}{cm^3}$, d'où
$m=9{,}6\times7{,}14=\res{\qty{69}{g}}$.
\item $n(\ce{Zn})=\dfrac{68{,}5}{65{,}4}=\qty{1,05}{mol}$, $n(\mathrm{e}^-)=\qty{2,10}{mol}$,
$Q=\qty{2,02e5}{C}$ et $\Delta t=\dfrac{\num{2.02e5}}{200}=\qty{1,0e3}{s}$, soit \res{\text{environ 17~min}}.
\item Même rayé, le revêtement protège : le zinc, \res{\text{plus réducteur}}, est oxydé à la place du fer (anode
sacrificielle). La rouille n'apparaît que lorsque le zinc est épuisé.
\end{questions}
\end{corrige}
\end{blocLivre}
\begin{blocLivre}{cor-10-12}
\begin{corrige}{12}{Recharger une voiture électrique}
\demarche{Énergie d'une cellule : $E=U\,Q$ ; la charge insère $n(\mathrm{e}^-)=Q/F$ mole de lithium dans le graphite.}
\begin{questions}
\item $E=U\,Q=3{,}7\times5{,}0=\res{\qty{18,5}{W.h}}$ ; nombre de cellules :
$\dfrac{\num{60000}}{18{,}5}\approx\res{\num{3,2e3}}$.
\item $\Delta t=\dfrac{60}{11}=\res{\qty{5,5}{h}}$ environ.
\item Il s'y produit la réduction \ce{Li+ + e- + 6 C -> LiC6} : le graphite est la \res{\text{cathode}}, reliée à la
borne \res{-} du chargeur.
\item $n(\mathrm{e}^-)=\dfrac{5{,}0\times3600}{\num{96500}}=\qty{0,187}{mol}$ : il s'insère autant d'ions lithium,
soit $m=0{,}187\times6{,}94=\res{\qty{1,3}{g}}$ de lithium par cellule.
\item La même réaction se produit alternativement dans les deux sens, des milliers de fois : c'est cette
\res{\text{réversibilité}} qui permet de stocker puis de restituer l'énergie, cycle après cycle.
\end{questions}
\end{corrige}
\end{blocLivre}
\begin{blocLivre}{cor-10-13}
\begin{corrige}{13}{Le dihydrogène vert}
\demarche{Chaque cellule en série produit I/(2F) mole de dihydrogène par seconde.}
\begin{questions}
\item En combinant les deux demi-équations (×2 pour la cathode) : \ce{2 H2O -> 2 H2 + O2}. Puissance :
$P=160\times2{,}0\times\num{25000}=\res{\qty{8,0}{MW}}$.
\item Par cellule : $\dfrac{I}{2F}=\dfrac{\num{25000}}{\num{193000}}=\qty{0,130}{mol.s^{-1}}$ ; pour 160 cellules :
$\res{\qty{20,7}{mol.s^{-1}}}$.
\item Par jour : $n=20{,}7\times86400=\qty{1,79e6}{mol}$, soit $m=\num{1.79e6}\times2{,}0=\res{\qty{3,6}{t}}$.
\item Énergie par jour : $E=8{,}0\times24=\qty{192}{MW.h}$, soit $\dfrac{\num{192000}}{3581}=\res{\qty{54}{kW.h}}$ par
kilogramme. La combustion n'en restitue que 33,3 : le rendement du stockage vaut
$\dfrac{33{,}3}{53{,}6}=\res{62~\%}$ ; le reste est dissipé en chaleur.
\item Il faut que l'électricité consommée soit \res{\text{d'origine renouvelable}} (ou décarbonée) : l'électrolyse ne
fait que stocker l'énergie électrique qu'on lui fournit.
\end{questions}
\end{corrige}
\end{blocLivre}
\begin{blocLivre}{cor-10-14}
\begin{corrige}{14}{De l'aluminium par électrolyse}
\demarche{Trois électrons par atome d'aluminium, quatre par molécule de CO2 à l'anode.}
\begin{questions}
\item $\res{\ce{Al^3+ + 3 e- -> Al}}$.
\item $n(\mathrm{e}^-)=\dfrac{\num{350000}\times86400}{\num{96500}}=\qty{3,13e5}{mol}$ ;
$n(\ce{Al})=\qty{1,04e5}{mol}$ et $m=\num{1.04e5}\times27{,}0=\res{\qty{2,8}{t}}$ par jour.
\item $n(\ce{C})=n(\ce{CO2})=\dfrac{n(\mathrm{e}^-)}{4}=\qty{7,8e4}{mol}$ : $m(\ce{C})=\res{\qty{940}{kg}}$ et
$m(\ce{CO2})=\num{7.8e4}\times44{,}0=\res{\qty{3,4}{t}}$ par jour.
\item $E=U\,I\,\Delta t=4{,}2\times\num{350000}\times86400=\qty{1,27e11}{J}=\qty{35,3}{MW.h}$ par jour, soit
$\dfrac{35{,}3}{2{,}82}=\res{\qty{12,5}{MW.h}}$ par tonne d'aluminium.
\item Le recyclage évite 95~\% de cette énergie et les émissions directes de \ce{CO2} des anodes : refondre
l'aluminium usagé est \res{\text{vingt fois moins énergivore}} que l'électrolyse.
\end{questions}
\end{corrige}
\end{blocLivre}
\begin{blocLivre}{cor-10-15}
\begin{corrige}{15}{Stocker l'énergie d'une maison autonome}
\demarche{Le rendement d'une chaîne est le produit des rendements ; la masse de dihydrogène découle de l'énergie stockée.}
\begin{questions}
\item Batteries : restitution $720\times0{,}90=\res{\qty{648}{kW.h}}$ ; masse
$\dfrac{\num{720000}}{150}=\qty{4800}{kg}$, soit \res{\qty{4,8}{t}} de batteries.
\item Rendement global : $0{,}70\times0{,}50=\res{35~\%}$ ; restitution $720\times0{,}35=\res{\qty{252}{kW.h}}$.
\item L'électrolyseur convertit $720\times0{,}70=\qty{504}{kW.h}$ en énergie chimique :
$m=\dfrac{504}{33{,}3}=\res{\qty{15,1}{kg}}$ de dihydrogène, soit $n=\dfrac{\num{15100}}{2{,}0}=\qty{7,6e3}{mol}$.
\item $V=\dfrac{nRT}{P}=\dfrac{7560\times8{,}314\times298}{\num{3.0e7}}=\res{\qty{0,62}{m^3}}$ : un réservoir de
620~L environ.
\item Les batteries restituent plus d'énergie, mais 4,8~t de batteries immobilisées des mois s'autodéchargent et
vieillissent : pour un stockage \res{\text{saisonnier}}, le dihydrogène conserve l'énergie indéfiniment dans un volume
raisonnable. C'est ce critère de durée, absent des calculs, qui le favorise.
\end{questions}
\end{corrige}
\end{blocLivre}
